
\subsection{Constraint-Aware Diffusion Models for Trajectory
Optimization}
\subsubsection{Abstract}
Abstract. The diffusion model has shown success in generating highquality and diverse solutions to trajectory optimization problems. However, diffusion models with neural networks inevitably make prediction
errors, which leads to constraint violations such as unmet goals or collisions. This paper presents a novel constraint-aware diffusion model for
trajectory optimization. We introduce a novel hybrid loss function for
training that minimizes the constraint violation of diffusion samples compared to the groundtruth while recovering the original data distribution.
Our model is demonstrated on tabletop manipulation and two-car reachavoid problems, outperforming traditional diffusion models in minimizing
constraint violations while generating samples close to locally optimal
solutions.
\subsubsection{Interesting Points}
\begin{itemize}
    \item Aim of this paper is to solve the problem of constraint violation. This is done by introducing a new loss (not only denoising loss) that learns constraint violation.
    \item To not hinder the process of inference however (where early, very noise states are naturally constraint violated), the loss is always calculated w.r.t. the level of noise that is present in the current step.
    \item This probably does not go well hand in hand with the progressive distillation?
\end{itemize}